\section{Historical Sources and Concordance}

The six questions used throughout this guide are a modern pedagogical discipline, never a scheme lifted from Aristotle or St.\ Thomas Aquinas. Whenever a modern reading term stands beside a classical term of art, the difference column prevents projecting modern habits backward into ancient and medieval texts.

\subsection{Concordance: Reading Terms and Historical Antecedents}

\begin{fourstable}{Reading Term}{Historical Antecedent}{Overlap}{Difference}
Warrant & Six conditions on demonstrative premises and kinds of principle (\emph{Post.\ An.}\ I.2, 71b17--72a25) & Both ask how a premise is known rather than merely asserted. & Aristotle defines a finished science; our eight kinds classify what an author actually offers on the page, including revelation and provenance only. \\
Burden & Demand for primary, immediate premises refusing circle and regress (\emph{Post.\ An.}\ I.2--3, 71b17--73a20) & Both hunt the proposition on which everything else rests. & Aristotle states a condition on completed demonstration; our Burden diagnoses a page that may be no demonstration at all. \\
Consequence & Commensurate universal (\emph{Post.\ An.}\ I.4, 73b27--74a3); each \emph{via} concluding \emph{that}, not \emph{what} (\emph{ST} I, q.\,2, a.\,3) & Both hold a conclusion strictly to the scope earned. & Conspicuous \emph{Not yet established} is a reader's check; classical authors leave attributes to separate questions. \\
Denial & Objection and reply, \emph{disputatio}, elenchus of non-contradiction (\emph{Metaph.}\ Gamma 4, 1006a28ff.) & Both measure a position by what rejecting it costs. & \emph{Disputatio} is a live exchange; our five kinds are a private classification applied to a page. \\
\emph{Quia} / \emph{propter quid} & \emph{Demonstratio quia} and \emph{propter quid} (\emph{Post.\ An.}\ I.13, 78a22--79a16; \emph{ST} I, q.\,2, a.\,2) & Taken directly from the tradition rather than invented. & Here the pair grades individual passages; Aristotle classifies modes within a science, and Aquinas licenses effect-to-cause proof without grasping essence. \\
Demonstration & \emph{Apodeixis} / \emph{demonstratio} (\emph{Post.\ An.}\ I.2, 71b18; \emph{ST} I, q.\,2, a.\,2) & Both mean causal knowledge, not mere validity. & Classical demonstration requires necessity, immediacy, and a \emph{per se} middle; modern speech often means only valid deduction. \\
Principle & \emph{Arche} / \emph{principium} (\emph{Post.\ An.}\ I.2, 72a8) & Both name an undemonstrated starting point. & Aristotle's principles are immediate and necessary; our ``metaphysical principle'' covers any substantive starting claim, contested or not. \\
Dialectic & Reasoning from \emph{endoxa} (\emph{Topics} I.1, 100a18--b23) & Both name argument from what will be granted rather than causes. & For Aristotle dialectic is a trained art; here it names a status a passage may have regardless of genre. \\
\end{fourstable}

\textbf{Claim}, \textbf{Meaning}, the plain-language reading, and \textbf{Not yet established} have no direct historical antecedent as frame divisions, though Aristotle demands fixing what a name signifies before arguing (\emph{Post.\ An.}\ I.1, 71a11--17).

\subsection{Verified Primary Passages Used in the Lessons}

Bekker page and line numbers govern Aristotle (Oxford translations by G.\ R.\ G.\ Mure, W.\ D.\ Ross, and W.\ A.\ Pickard-Cambridge). Aquinas follows the Leonine Latin editions (\emph{In PA} 1882; \emph{ST} I 1888; \emph{SCG} 1961), Baur/Koch (\emph{De ente} 1933), Decker 1959 (\emph{Super Boet.\ De Trin.}), and the English Dominican translation of \emph{ST} I unless explicitly marked as a close rendering preserving Latin clause structure.

\begin{threestable}{Author and Passage}{Verified Doctrine}{Reading Caution}
Aristotle, \emph{Post.\ An.}\ I.1--4, 71a11--74a3 & Twofold prior knowledge (71a11--17); \emph{episteme} and \emph{apodeixis} from true, primary, immediate premises better known in nature, prior to, and causes of the conclusion (71b8--72a25); no circle or regress (72b5--73a20); \emph{per se} predication and commensurate universal (73a34--74a3). & The threefold \emph{an sit} / \emph{quid sit} / \emph{quia est} triad is later scholastic formulation. \emph{Per se} predication never licenses a \emph{per se} causal series. \\
Aristotle, \emph{Post.\ An.}\ I.6--8, 13--14, 30; II.8, 11, 16, 19 & Necessary premises and single subject genus barring subalternation (74b5--75b36, 87b19--27); \emph{quia} vs.\ \emph{propter quid} via planets/twinkling (78a30--38) and moon/spherical (78b3--10); first figure (79a17--32); four causes as middle (94a20--95a9); deciduous/broad leaves reciprocal reversal (II.16, 98b4--10); \emph{epagoge} and \emph{nous} (99b18--100b17; \emph{Nic.\ Eth.}\ VI.6). & Deciduousness/vine stands in Book II (98b4--10), not I.13. The reading of II.19 (\emph{epagoge}/\emph{nous}) remains genuinely disputed among scholars. \\
Aristotle, \emph{Topics} I.1--2; \emph{Metaph.}\ Gamma 1--4 & Dialectic argues from \emph{endoxa} and examines first principles (100a18--b23, 101a25--37); non-contradiction defended elenctically (1005b15--1009a5). & \emph{Probabilia} means reputable, not numerically probable. \\
Aquinas, \emph{In PA} I--II; \emph{Super Boet.\ De Trin.}\ q.\,5 a.\,3, q.\,6 a.\,4 ad 2 & Literal commentary (\emph{lectiones}); cause as middle (\emph{In PA} II, lect.\ 1 n.\,8); metaphysics (\emph{ens commune}) attained by negative judgment (\emph{separatio}), distinct from two modes of abstraction; effect-to-cause proof without essence (\emph{De Trin.}\ q.\,6 a.\,4 ad 2). & ``Three degrees of abstraction'' and \emph{via inventionis}/\emph{via iudicii} are later manual labels, not Aquinas's wording here. \\
Aquinas, \emph{ST} I, q.\,1 aa.\,2, 8 ad 2; q.\,2 aa.\,1--3; qq.\,3--11, 14 & Sacred doctrine from higher science; philosophers probable, Scripture necessary; \emph{per se notum quoad se} vs.\ \emph{quoad nos} (Obj.\ 1 Damascene; \emph{sed contra} Ps.\ 52:1; ad 2 refuses name-to-reality leap); \emph{duplex est demonstratio}, \emph{praeambula fidei}, \emph{quid significet nomen} as middle; Five Ways (\emph{viae}); entitative attributes in qq.\,3--11, intellect in q.\,14. & Aquinas does not use \emph{subalterna} at q.\,1 a.\,2, does not name Anselm at q.\,2 a.\,1, and calls his arguments \emph{viae} rather than ``Five Proofs.'' \\
Aquinas, \emph{ST} I, q.\,13 aa.\,3, 5; q.\,44 aa.\,1--2; q.\,46 a.\,2; \emph{SCG} I cc.\,3, 13; \emph{De ente} c.\,4 & Analogy midway between equivocation and univocity; participation argument; world's temporal beginning \emph{sola fide}; \emph{per se} vs.\ \emph{per accidens} efficient regress (q.\,46 a.\,2 ad 7); ordered/instrumental movers and efficient causes (\emph{SCG} I c.\,13 nn.\,14, 15, 33); essence/existence dependence (\emph{De ente} c.\,4). & The syllogistic four-term problem for analogical middles belongs to Cajetan (\emph{De nominum analogia}, 1498), not q.\,13. Unchecked texts (\emph{In Phys.}\ VIII, \emph{SCG} II c.\,38) are excluded. \\
Bl.\ John Duns Scotus, \emph{Ord.}\ I d.\,3 p.\,1 qq.\,1--2 n.\,26; David Hume, \emph{Dialogues}, Part IX (1779) & Scotus defines univocity by sufficiency for contradiction and a syllogistic middle; Hume's Cleanthes argues explaining each of twenty particles explains the whole. & Non-Thomistic comparators in Lesson 7 comparing burdens across frameworks. \\
\end{threestable}
